\section{Introduction}\label{sec:intro}

A \emph{magma} is a set equipped with a single binary operation subject to no
axioms.  Despite this minimal structure, the lattice of equational theories on
magmas is remarkably rich: the Equational Theories Project \cite{ETP2025}
completely determined the $4{,}694 \times 4{,}694$ implication matrix for all
equational laws of order at most four, formalizing every proof in Lean.  Of the
$22{,}033{,}636$ ordered pairs (including reflexive), $8{,}178{,}279$ (37.1\%)
are true implications and $13{,}855{,}357$ (62.9\%) are false.

This complete determination, while a tour de force of collaborative formal
mathematics, produces a matrix that is opaque to human understanding.  The
natural question arises: \emph{to what extent can the implication structure be
explained by simple structural properties of the equations themselves?}

This question has both theoretical and practical motivation.  Theoretically,
understanding which syntactic features of an equation determine its
implications sheds light on the algebraic content encoded in equational laws.
Practically, the Mathematics Distillation Challenge \cite{Tao2026blog,SAIR2026}
asks whether the implication structure can be compressed into a short text
(under 10KB) that enables even a weak language model to predict implications
--- a test of whether mathematical knowledge admits compact, interpretable
representations.

\subsection{Prior work}
The Equational Theories Project \cite{ETP2025} established the complete
implication graph using a combination of techniques: finite magma
counterexamples (refuting 96.3\% of false implications with only 524 distinct
magmas of size $\le 4$), syntactic rewriting, matching invariants, automated
theorem provers (Vampire, Prover9, Z3), equational reasoning via e-graphs, and
SAT solvers.  They showed that the $42$\,MB bit table compresses to $40$\,kB
with LZMA2, and that a five-layer convolutional neural network achieves 99.7\%
prediction accuracy from equation syntax alone, demonstrating high
compressibility.

Berlioz and Melli\`{e}s \cite{BM2026} introduced a \emph{Stone pairing}
approach, computing for each equation the probability that a random variable
assignment satisfies it in a randomly sampled finite magma.  PCA on the
resulting probability matrix produces a three-dimensional latent space in which
equivalent theories cluster tightly and implications flow in a structured,
oriented manner.

Aliniaeifard and Li \cite{AL2025} developed neural methods for identifying
sufficient conditions for mathematical implications, while Olausson et
al.~\cite{LINC2023} showed that combining language model parsing with external
logical reasoning outperforms pure neural approaches on logical inference tasks.

\subsection{Our contributions}
We systematically extract decision rules from the full implication matrix and
prove their correctness.  Our contributions are as follows.

\begin{itemize}[leftmargin=*,itemsep=2pt,topsep=2pt]
\item We introduce a five-class \emph{structural form} classification of
  equational laws (Definition~\ref{def:form}) and prove that form-based
  implication rules cover 39.8\% of all non-reflexive pairs with zero error
  (Theorem~\ref{thm:absorbing}, Theorem~\ref{thm:general-nonimplication}).

\item We establish a \emph{signature direction principle}
  (Theorem~\ref{thm:sig-direction}) showing that equations with operations on
  the left-hand side never imply equations with a bare variable on the
  left-hand side, covering an additional 22\% of pairs.

\item We prove that the difference in distinct variable counts is the dominant
  heuristic predictor for the remaining cases
  (Proposition~\ref{prop:nvar-monotone}) and provide calibrated probability
  estimates.

\item We compose these rules into an ordered decision procedure that achieves
  88.9\% accuracy on random samples, compared to 62.9\% for the always-false
  baseline.
\end{itemize}

\subsection{Paper organization}
Section~\ref{sec:prelim} introduces the necessary definitions and notation.
Section~\ref{sec:main} presents our main results: the structural classification,
form-based implication rules, signature direction principle, variable count
analysis, and the combined decision procedure with its empirical evaluation.
Section~\ref{sec:discussion} discusses limitations, connections to related
problems, and open questions.  Section~\ref{sec:conclusion} summarizes our
findings.
